\newglossaryentry{api_rest}{
    name={REST API},
    description={Application Programming Interface based on the \textit{Representational State Transfer} (REST) paradigm, which uses the HTTP protocol to facilitate communication between systems and allows performing operations such as querying, creating, modifying and deleting resources in a scalable and efficient manner}
}

\newglossaryentry{gitops}{
    name={GitOps},
    description={Operational paradigm that uses Git repositories as the single source of truth for defining and managing the desired state of infrastructure and applications. Any change to the system is made through commits in the repository, and an automated operator (such as ArgoCD) is responsible for synchronising the actual state of the environment with the state declared in Git}
}

\newglossaryentry{cni}{
    name={CNI},
    description={(\textit{Container Network Interface}) Specification and set of libraries that define how network plugins should configure network interfaces of containers in Kubernetes. Examples of CNI plugins are Calico, Flannel and Cilium}
}

\newglossaryentry{drift_protection}{
    name={Drift protection},
    description={Mechanism that automatically detects and corrects any deviation between the desired state of a resource (defined declaratively) and its actual state in the system. In the context of Kyverno, it is implemented through the \texttt{synchronize: true} field in \textit{generate} policies, which automatically regenerates resources that are deleted or modified in an unauthorised manner}
}

\newglossaryentry{auto_hardening}{
    name={Auto-hardening},
    description={Automatic process of applying security configurations to system resources without manual intervention from the operator or developer. In the context of this work, it is implemented through Kyverno mutation policies that inject security fields into Kubernetes manifests at deployment time}
}

\newglossaryentry{imds}{
    name={IMDS},
    description={(\textit{Instance Metadata Service}) Service available in cloud environments (AWS, GCP, Azure) accessible from any instance through the link-local IP address \texttt{169.254.169.254}. It exposes sensitive instance information such as temporary IAM credentials, service tokens and network configuration. It is a common attack vector in compromised container environments}
}

\newglossaryentry{supply_chain_attack}{
    name={Supply chain attack},
    description={Type of attack that compromises the software or tools used during the development or deployment process, rather than directly attacking the target system. In CI/CD environments, it can manifest as the manipulation of dependencies, container images, GitHub Actions or third-party runners to inject malicious code into the pipeline}
}

\newglossaryentry{policy_report}{
    name={PolicyReport},
    description={Native Kubernetes resource automatically generated by Kyverno that consolidates the results of policy evaluation over cluster resources. It indicates which resources have passed, failed or been mutated by each policy, with a timestamp for each evaluation. There are two variants: \texttt{PolicyReport} (namespace-scoped) and \texttt{ClusterPolicyReport} (cluster-scoped)}
}

\newglossaryentry{webhook}{
    name={Webhook},
    description={Kubernetes API Server extension mechanism that allows invoking external HTTP services during the resource admission process. In the context of Kyverno, admission webhooks are the channel through which the API Server delegates policy evaluation to the Kyverno engine. There are two types: \texttt{MutatingWebhookConfiguration} (for mutation policies) and \texttt{ValidatingWebhookConfiguration} (for validation policies)}
}

\newglossaryentry{etcd}{
    name={etcd},
    description={Distributed key-value database that acts as the persistent store for the complete Kubernetes cluster state. All cluster objects (pods, services, configmaps, policies, etc.) are persisted in etcd once the admission process has been passed. Its consistency and availability are critical for cluster operation}
}

\newglossaryentry{hardening}{
    name={Hardening},
    description={Process of reducing the attack surface of a system by eliminating or restrictively configuring unnecessary functionalities. In the context of Kubernetes containers, it includes measures such as disabling privilege escalation, removing unneeded Linux capabilities, enforcing execution as a non-root user and restricting filesystem access}
}

\newglossaryentry{cidr}{
    name={CIDR},
    description={(\textit{Classless Inter-Domain Routing}) Notation used to define IP address ranges. In Kubernetes, CIDR ranges are used to designate the pod address space (pod CIDR) and service address space (service CIDR). For example, \texttt{10.244.0.0/16} defines the range of IPs assignable to pods when using Calico as the CNI plugin}
}

\newglossaryentry{two_repo_gitops}{
    name={Two-repository pattern (GitOps)},
    description={GitOps architectural pattern that separates the application source code repository from the configuration repository containing Kubernetes manifests. The CI pipeline updates the configuration repository after a successful build, and ArgoCD synchronises the cluster with the state declared in that repository. This separation improves security by limiting cluster access permissions and clarifies the responsibilities of each component}
}

\newglossaryentry{qos_kubernetes}{
    name={QoS (Kubernetes)},
    description={(\textit{Quality of Service}) Classification that Kubernetes assigns to pods based on how their \texttt{resources.requests} and \texttt{resources.limits} are defined. There are three classes: \texttt{Guaranteed} (requests == limits for all containers --- maximum protection against eviction), \texttt{Burstable} (requests < limits or only one defined) and \texttt{BestEffort} (no fields defined --- first class to be evicted under resource pressure)}
}

\newglossaryentry{patchstrategicmerge}{
    name={patchStrategicMerge},
    description={Kyverno patch application mechanism that performs an intelligent merge (\textit{merge}) of the original Kubernetes object with the patch defined in the policy. Unlike a complete replacement, \texttt{patchStrategicMerge} only adds absent fields and respects values already explicitly declared in the original manifest, guaranteeing that mutations are additive and non-destructive}
}

\newglossaryentry{runner}{
    name={Runner (GitHub Actions)},
    description={Execution agent that processes jobs defined in GitHub Actions workflows. It can be a GitHub-hosted runner (\textit{hosted runner}) or a self-managed runner on own infrastructure (\textit{self-hosted runner}). Runners execute each pipeline step in an isolated environment and are the target of supply chain attacks that seek to compromise the build process or exfiltrate CI environment credentials}
}

\newglossaryentry{multi_tenant}{
    name={Multi-tenant},
    description={Kubernetes deployment model in which a single cluster is shared by multiple teams, applications or clients (called \textit{tenants}) with logical isolation through namespaces, RBAC and Network Policies. This model optimises resource usage but introduces security challenges related to tenant isolation, especially in scenarios where a workload compromise could propagate to others}
}
